% █▀▀ █▀   █▀█ █▀█ █▀▀ ▄▀█ █▀▄▀█ █▄▄ █   █▀▀
% █▄▄ ▄█   █▀▀ █▀▄ ██▄ █▀█ █ ▀ █ █▄█ █▄▄ ██▄

% NOTE: The packages (and related commands) in this preamble are loaded following the bellow ordering.
% Packages related to:
	% 1. Document canvas        >      i.e.  page geometry            
	% 2. Typeface system        >      i.e.  used fonts               
	% 3. Document structure     >      i.e.  headers formatting       
	% 4. Floats                 >      i.e.  non-textual elements     
	% 5. Micro-Layout           >      i.e. "inline" adjustments     
	% 6. Finalization           >      i.e.  packages best loaded last

%::::::::::::::::::::::::::::::::::::::::::::::::::::::::::::::::::::::::::::::::
%                                   1. CANVAS                                    
%           Define document's physical space plus some basic utilities           
%................................................................................
% CANVAS
\usepackage{geometry}		% Page geometry (margins configured in main file)
\usepackage[brazil]{babel}	% Adjust document to the language it's written in

% BASIC UTILITIES
\usepackage{csquotes}		% Language-sensitive quoting;
\usepackage[dvipsnames]{xcolor}	% Access to 68 named colors
\usepackage{moresize}		% Additional font sizes
\usepackage[normalem]{ulem}	% More underline styles (and aligns bottoms of multiple underlines)
				% 'normalem' prevents underlines from altering \emph


%::::::::::::::::::::::::::::::::::::::::::::::::::::::::::::::::::::::::::::::::
%                               2. TYPEFACE SYSTEM                               
%    Define all fonts - for text, algorithms and math - used in the document     
%................................................................................
% 2.1 SYSTEM BACKEND
\usepackage{fontspec}		% Native Unicode-to-OpenType mapping;
				% Replaces legacy TeX font-enconding hacks;
				% DOCUMENT MUST NOW BE COMPILED WITH Lua OR XeLaTex

\usepackage{amsmath}		% AMS's package for mathematical typesetting
\usepackage{mathtools}		% Extends amsmath's functionalities (i.e.: a superset)
\usepackage[
	math-style = ISO,
	bold-style = ISO
]{unicode-math}			% Must be loaded after 'mathtools';
				% DOCUMENT MUST NOW BE COMPILED WITH Lua OR XeLaTex

%. . . . . . . . . . . . . . . . . . . . . . . . . . . . . . . . . . . . . . . . . 
% 2.2 SYSTEM DEFINITION
% 2.2.1 TEXT FONTS
\setmainfont{TeX Gyre Pagella}[Scale = 1.0]
\setsansfont{URW Classico}[
	Scale = MatchLowercase,	% Ensure main and sans fonts are at the same scale
]
\setmonofont{Inconsolata}[Scale=MatchLowercase]

% 2.2.2 MATH FONTS
\setmathfont{Libertinus Math}[Scale=MatchLowercase]
% Manually pick distinct calligraphic (Chancery hand) and script (Round hand) fonts for a greater variety of symbols
\setmathfont{XITS Math}[
	Scale 		= MatchLowercase,
	range		= {cal,bfcal},
	StylisticSet	= 1
]
\setmathfont{Libertinus Math}[
	Scale 		= MatchLowercase,
	range		= {scr,bfscr}
]
% Define a custom font face 'mathopfont' and use it to distinguish mathematical operators from prose
\setmathfontface\mathopfont{Latin Modern Roman}[Scale=MatchLowercase]
\setoperatorfont\mathopfont


%::::::::::::::::::::::::::::::::::::::::::::::::::::::::::::::::::::::::::::::::
%                             3. DOCUMENT STRUCTURE                              
%         Article structure assumed. Here, we just stylise their headers         
%................................................................................
\usepackage{titlesec}		% Controls vertical rhythm and heading aesthetics

% STYLISE HEADINGS
\titleformat{\section}
    {\sffamily\bfseries}	% Format
    {\thesection}		% Label
    {1em}			% Space between number and title
    {\MakeUppercase}		% Before-code
\titleformat{\subsection}
    {\sffamily\bfseries\small}	% Format
    {\thesubsection}    	% Label
    {1em}               	% Space between number and title
    {\MakeUppercase}    	% Before-code
%\titleformat{\subsubsection}
%    {\sffamily\bfseries\small}	% Format
%    {\thesubsubsection}	% Label
%    {1em}               	% Space between number and title
%    {\MakeUppercase}    	% Before-code
%\titleformat{\paragraph}[runin]
%    {\sffamily\bfseries}	% Format
%    {\theparagraph}		% Label
%    {1em}			% Space between number and title
%    {}				% Before-code
%\titlespacing*{\paragraph}
%    {0pt}			% Left margin indentation
%    {3.25ex plus 1ex minus .2ex}% Space BEFORE the paragraph title (vertical)
%    {1em}			% Space AFTER the title (horizontal)

%. . . . . . . . . . . . . . . . . . . . . . . . . . . . . . . . . . . . . . . . . 
% PAGINATIONG AND TRACKING
\usepackage{fancyhdr}	% Custom headers and footers
\pagestyle{fancy}
\fancyhf{}				% Clear all default headers/footers
% Defensive practice to keep my pages clearly identifiable
% if the funding agencies' dossier gets scrambled
\renewcommand{\headrulewidth}{0pt}	% Remove the default top line
\fancyfoot[C]{\sffamily\small \thepage}
\fancyfoot[R]{\sffamily\small Cody S. B. Setti -- Súmula Curricular}


%::::::::::::::::::::::::::::::::::::::::::::::::::::::::::::::::::::::::::::::::
%                                    4. FLOATS                                   
%               Manage Floating Content (e.g.: figures and tables)               
%................................................................................
% FLOATS IN GENERAL
\usepackage{float}	% Provides [H] 'Here' float placement specifier (use with caution)
\usepackage{subcaption}	% Format multiple images (tables) in one figure (table) environment
			% NOTE: This package *automatically* loads the 'caption' package,
\captionsetup{
	labelfont = {small,bf},
	textfont  = {small,it}
}
%\subcaptionsetup{...}	% Uncomment to customize sub-captions separately

%. . . . . . . . . . . . . . . . . . . . . . . . . . . . . . . . . . . . . . . . . 
% TABLES
\usepackage{booktabs}	% Format tables following scientific convention
\usepackage{multirow}	% Allow multirow cells in tables
\usepackage{tabularx}	% Auto-adjust column widths to fill a specified amount


%::::::::::::::::::::::::::::::::::::::::::::::::::::::::::::::::::::::::::::::::
%                             5. MICRO-LAYOUT                              
%                       Packages that alter "inline" text                        
%................................................................................
\usepackage{enumitem}		% Customize list environments (provides 'resume')
\setlist{nosep, leftmargin=*}	% Compresses list vertical spacing globally

\usepackage{multicol}		% Environment with multiple columns of text
\usepackage{fontawesome5}	% For Orcid's logo
\usepackage{academicons}	% For Google Scholar's logo
\usepackage{graphics}		% For manually inserting Lattes' and WoS' logos


%:::::::::::::::::::::::::::::::::::::::::::::::::::::::::::::::::::::::::::::::::
%                                 6. FINALIZATION                                 
%                 Packages that patch others. MUST BE LOADED LAST                 
%.................................................................................
% BIBLIOGRAPHY                                   
\usepackage[
	backend	 = biber,	% Use Biber backend (more powerful than BibTex)
	style	 = numeric,
	backref	 = true,
	maxnames = 15,		% Enforces the FAPESP 15-author maximum rule
]{biblatex}

%. . . . . . . . . . . . . . . . . . . . . . . . . . . . . . . . . . . . . . . . .
% HYPERLINKS
\usepackage{hyperref}	% Add hyperlinks to references, ToC and URLs
\hypersetup{
	colorlinks	= true,		% true: colored links; false: boxed links
	linkcolor	= NavyBlue,	% Internal links
	citecolor	= NavyBlue,
	urlcolor	= NavyBlue,
	breaklinks	= true,
	pdfborder	= {0 0 0}
}

%. . . . . . . . . . . . . . . . . . . . . . . . . . . . . . . . . . . . . . . . .
% MICRO-TYPOGRAPHY
\usepackage{microtype}		% Enhance micro-typography (e.g.: protrusion, expansion)
				% Loaded last for best compatibility with other packages
